\chapter{Grundlagen}

\todox[]{Intro schreiben}

\section{ImageNet als Innovationstreiber in der Bildklassifikation}

Als eines der bekanntesten Daten- und Benchmarking-Quellen in der Bildklassifikation gilt ImageNet, ein Projekt der Universitäten Stanford und Princeton \cite{imagenet_homepage}. Im Rahmen dieser Initiative wurden über die Jahre hinweg mehrere Datensätze veröffentlicht, auf denen Forscher ihre Modelle trainiert und getestet haben. Beispielsweise wurden 2009 zirka 3.2 Millionen Bilder zur Verfügung gestellt, die 5427 Kategorien zugeordnet wurden \cite{deng2009imagenet}.

Zwischen 2012 und 2017 wurden die sogennanten ImageNet Large Scale Visual Recognition Challenge (ILSVRC) durchgeführt, bei der Forscher ihre Modelle auf einem standardisierten Datensatz mit \enquote{ 1,281,167 training images, 50,000 validation images and 100,000 test images} \cite{imagenet_download} testen und vergleichen konnten \cite{imagenet_lsvrc_challenge, imagenet_homepage}. Diese Challenges haben zu bedeutenden Fortschritten in der Bildklassifikation geführt: Forscher haben Daten erhalten und eine standardisierte Umgebung, in der sie ihre Techniken testen und vergleichen konnten, was zu einem intensiven Wettbewerb geführt hat, in dem viele innovative Architekturen0 entsanden sind.
Zwischen 2012 und 2017 wurden die sogennanten ImageNet Large Scale Visual Recognition Challenge (ILSVRC) durchgeführt, bei der Forscher ihre Modelle auf einem standardisierten Datensatz mit \enquote{ 1,281,167 training images, 50,000 validation images and 100,000 test images} \cite{imagenet_download} testen und vergleichen konnten \cite{imagenet_lsvrc_challenge, imagenet_homepage}. Diese Challenges haben zu bedeutenden Fortschritten in der Bildklassifikation geführt: Forscher haben Daten erhalten und eine standardisierte Umgebung, in der sie ihre Techniken testen und vergleichen konnten, was zu einem intensiven Wettbewerb geführt hat, in dem viele innovative Architekturen0 entsanden sind.
Zwischen 2012 und 2017 wurden die sogennanten ImageNet Large Scale Visual Recognition Challenge (ILSVRC) durchgeführt, bei der Forscher ihre Modelle auf einem standardisierten Datensatz mit \enquote{ 1,281,167 training images, 50,000 validation images and 100,000 test images} \cite{imagenet_download} testen und vergleichen konnten \cite{imagenet_lsvrc_challenge, imagenet_homepage}. Diese Challenges haben zu bedeutenden Fortschritten in der Bildklassifikation geführt: Forscher haben Daten erhalten und eine standardisierte Umgebung, in der sie ihre Techniken testen und vergleichen konnten, was zu einem intensiven Wettbewerb geführt hat, in dem viele innovative Architekturen und Verfahren entstanden sind.


\section{Architekturen für die Bildklassifikation}
Mit der zunehmenden Verfügbarkeit von großen Bilddatensätzen und der Entwicklung leistungsfähiger Hardware haben sich Deep Learning-Modelle als äußerst effektiv für die Bildklassifikation erwiesen \cite{alex2012alexnet}. In den folgenden Abschnitten werden einige der wichtigsten Architekturen vorgestellt, die in diesem Bereich verwendet werden.

\subsection{Convolutional Networks}
Im Vergleich zu den klassischen \textit{Feedfoward Neural Networks} (FNNs) sind \textit{Convolutional Neural Networks} (CNNs) effizienter: Sie haben eine effiziente Struktur mit weniger Gewichten und Kanten zwischen den Modell-Schichten etabliert, was das Training erleichtert, ohne dabei einen betrachtlichen Verlust an Leistung zu verursachen. \cite{alex2012alexnet}
Dadurch eignen sich CNNs besonders gut für die Verarbeitung von Bilddaten, da sie die räumlichen Beziehungen zwischen den Pixeln nutzen können, um relevante Merkmale durch verschiedene Abstraktionsebenen zu extrahieren. 

Das sogenannte AlexNet \cite{alex2012alexnet} war eines der ersten CNNs, das auf dem ImageNet-Datensatz eine signifikante Verbesserung der Bildklassifikationsleistung erzielte und damit den Durchbruch für CNNs in diesem Bereich markierte. Seitdem wurden viele weitere Architekturen entwickelt, die sich von AlexNet bei der Vewendung von Convolutional-Layers inspiriert lassen haben, wie z. B. VGGNet \cite{simonyan2015vgg}, ResNet \cite{he2016resnet}, MobileNet \cite{howard2017mobilenet} und EfficientNetV2 \cite{tan2021efficientnetv2}.

Die nachsten Abschnitte konzentrieren sich auf die Architekturen ResNet und EfficientNetV2, da sie aufgrund ihrer historischen und technischen Relevanz hohe Bedeutung für die Bildklassifikation haben.

\subsection{Residual Networks}
Die Erhöhung der Anzahl der Schichten in einem Deep Neural Network (DNN) war nicht immer mit einer Verbesserung der Vorhersagegüte verbunden.
In dieser Hinsicht wurde beobachtet, dass die Erweiterung eines Netzwerks um eine größere Anzahl von Schichten oft zu einer deutlichen Erhöhung der Fehlerrate sowohl beim Training vor als auch bei den Tests führte. So wurde das Phänomen des sogennante \textit{Degradationsproblem} von \cite{he2016resnet} wie folgt beschrieben: 
\begin{quote}
    Driven by the significance of depth, a question arises: Is learning better networks as easy as stacking more layers? [...]
    
    When deeper networks are able to start converging, a degradation problem has been exposed: with the network depth increasing, accuracy gets saturated (which might be unsurprising) and then degrades rapidly. Unexpectedly, such degradation is not caused by overfitting, and adding more layers to a suitably deep model leads to higher training error, as reported in [10, 41] and thoroughly verified by our experiments. 
\end{quote}

Das hat die wissenschaftliche Gemeinschaft vor ein Rätsel gestellt: Falls die zusätzlichen Schichten nicht zu einer erhöhten Leistungen führen könnten, wären sie theoretisch trotzdem in der Lage die Identitätsfunktion zu approximieren, sodass die Leistung des erweiterten Modells zumindest nicht viel schlechter als die des ursprünglichen Modells sein könnte.
Dies war aber in der Praxis häufig nicht der Fall und die Leistung verschlechterte sich sogar deutlich.


So hat \cite{he2016resnet} dieses Problem durch die Einführung sogenannter Residual-Blöcke gelöst. Diese Blöcke enthalten eine direkte Verbindung zwischen der Eingabe und der Ausgabe, sodass die dazwischenliegenden Schichten lediglich die Differenz (d.h das \textbf{Residuum}) zwischen der Eingabe und der Ausgabe lernen müssen, wie die sogenannte \textbf{Residualfunktion} darstellt:
\begin{equation}
    y_k = x_k + \mathcal{F}(x_k, \mathcal{W}_k)
\end{equation}
Dabei bezeichnet $x_k$ die Eingabe eines Blocks, $\mathcal{F}(x_k, \mathcal{W}_k)$ die durch die Schichten parametrisierte Residualfunktion und $y_k$ die Ausgabe des Blocks. Dieses Vorgehen ermöglicht einem Residual-Block das \textit{Identity-Mapping} mit geringer Rechenkomplexität zu realisieren, indem die Eingabe unverändert zur Ausgabe addiert wird.

Dies ist besonders vorteilhaft, wenn die Ausgabe eines Blocks ähnlich der Eingabe ist, da die Schichten dann nur lernen müssen, die Unterschiede zu modellieren, anstatt die gesamte Funktion zu lernen.

ResNets haben in vielen Bildklassifikationswettbewerben hervorragende Ergebnisse erzielt und sind nach wie vor eine der beliebtesten Architekturen für die Bildklassifikation. Im originalen Paper \cite{he2016resnet} wurde ein 152-Schichte-Deep-Neural-Network vorgestellt, das auf dem ImageNet-Datensatz eine Top-5-Fehlerquote von nur 3,57\% erreichte.

\subsection{EfficientNetV2}

EfficientNet ist eine Familie von Convolutional Neural Networks, die eine bessere Leistung bei geringerer Rechenkomplexität erzielen. So hat sich das Modell EfficientNetV2, die Nachfolgerversion von EfficientNetV1, als besonders leistungsfähig erwiesen: Es verfügt über eine Top-1-Genauigkeit von 85,7\% auf dem ImageNet-Datensatz und gleichzeitig 3-9 Mal schnelleres Trainingsverfahren und 6-8x weniger Gewichtsparameter als der Großteil ähnlicher Neural Networks im Stand der Technik zu dem Zeitpunkt \cite{tan2021efficientnetv2}.

Diese Errungenschaft wurde vor allem durch die Einführung von Automated Machine Learning anhand Neural Architecture Search (NAS) zur Optimierung der Architektur von EfficientNetV2 und durch die Einführung von \textit{Fused-MBConv}-Blöcken erreicht, indem es die Suche nach der besten Kombination von Schichten und Hyperparametern zur Leistungsoptimierung automatisiert hat \cite{tan2021efficientnetv2}. Darüber hinaus hat die Verwendung von \textit{Progressive Learning} beim Training – zunächst mit kleineren Bildern und weniger Regularisierung und dann mit größeren Bildern und mehr Regularisierung – sich als besonders effektiv erwiesen, um die Trainingszeit zu verkürzen, ohne dabei die Korrektklassifizierungsrate des Modells zu beeinträchtigen \cite{tan2021efficientnetv2}.


\section{XAI-Verfahren}
\todox[]{kurz Erklären, was XAI ist, warum es wichtig ist, und welche Ziele es verfolgt (Vertrauen, Nachvollziehbarkeit etc.)}
\todox[]{wichtige XAI-Verfahren nennen und kurz erläutern (z.B. Grad-CAM, LIME etc.)}